% ============================================================
%  header.tex -- packages and layout, reused from
%  docs/Expose_Template/header.tex so this thesis matches the
%  department's Exposé/thesis look (fonts, colours, title-page
%  macros, listings style). Adjust freely.
% ============================================================
\usepackage[utf8]{inputenc}
\usepackage[ngerman,english]{babel}
\usepackage{siunitx}
% --- siunitx set-up ---------------------------------------------------------
% Every quantity in the thesis goes through \SI/\si/\num, so the number-unit
% spacing, the upright unit shape and the line breaking are decided here once
% instead of at each occurrence.  The four settings below reproduce the way
% the quantities were already written by hand, so switching to siunitx changes
% who does the typesetting, not what the page looks like.
%   separate-uncertainty       95.75 +- 2.20, not 95.75(220)
%   separate-uncertainty-units 95.75 +- 2.20 %, not (95.75 +- 2.20) %
%   per-mode                   m/s, not m s^{-1}
%   range-phrase               5--13 m/s, not "5 to 13 m/s"
\sisetup{
  separate-uncertainty        = true,
  separate-uncertainty-units  = single,
  per-mode                    = symbol,
  range-units                 = single,
  range-phrase                = {\text{--}},
}
% Units that are not SI but are what this detector measures in.  Declaring them
% keeps them out of the text: an ADU or a px is then written the same way as a
% millisecond, and siunitx sets it upright with the right space in front.
\DeclareSIUnit{\ADU}{ADU}      % analogue-to-digital unit, one count of the ADC
\DeclareSIUnit{\px}{px}        % pixel, as a length/area on the sensor
\DeclareSIUnit{\pixel}{pixel}  % spelled out; \mega\pixel gives "Mpixel"
\DeclareSIUnit{\fps}{fps}      % frames per second, as the camera reports it
% siunitx deprecates its own \bar (the name collides with the math accent
% \bar and with amsmath's \Bar), so the pressure unit is declared here under
% a name that is free.  Written \milli\barpressure it still prefixes.
\DeclareSIUnit{\barpressure}{bar}
\usepackage[acronym]{glossaries}
\usepackage{tabularx}
\usepackage{booktabs}
% \begin{tabular}\toprule ... \midrule ... \bottomrule
\usepackage[labelfont={sf,bf},textfont={sf}]{caption}
\captionsetup{singlelinecheck=false}
\usepackage{subcaption}
\usepackage{subfloat}
\usepackage{amsmath,bm}
\usepackage{physics}          % \bra, \ket, \braket, \dv, \pdv
\usepackage[T1]{fontenc}
\usepackage{graphicx}
\usepackage{xcolor}
% TU Wien colors
\definecolor{TUblue}{RGB}{0,102,153}
\definecolor{TUblue1}{RGB}{84,133,171}
\definecolor{TUblue2}{RGB}{114,173,213}
\definecolor{TUblue3}{RGB}{166,213,236}
\definecolor{TUblue4}{RGB}{223,242,253}

\definecolor{TUlightblue}{RGB}{143,190,229}
\definecolor{TUlighterblue}{RGB}{222,231,236}

\definecolor{TUgreen}{RGB}{0,126,113}
\definecolor{TUgreen1}{RGB}{106,170,165}
\definecolor{TUgreen2}{RGB}{162,198,194}
\definecolor{TUgreen3}{RGB}{233,241,240}

\definecolor{TUmagenta}{RGB}{186,70,130}
\definecolor{TUmagenta1}{RGB}{205,129,168}
\definecolor{TUmagenta2}{RGB}{223,175,202}
\definecolor{TUmagenta3}{RGB}{245,229,239}

\definecolor{TUorange}{RGB}{225,137,34}
\definecolor{TUorange1}{RGB}{238,180,115}
\definecolor{TUorange2}{RGB}{245,208,168}
\definecolor{TUorange3}{RGB}{253,239,225}

\definecolor{TUgray}{RGB}{100,99,99}
\definecolor{TUgray1}{RGB}{157,156,156}
\definecolor{TUgray2}{RGB}{208,208,208}
\definecolor{TUgray3}{RGB}{237,237,237}

\usepackage[export]{adjustbox}
\usepackage{placeins}

% Float placement. LaTeX's defaults are conservative enough that a figure with a
% long caption is often sent to a page of its own, with no text on it at all --
% which is where most of this thesis' white space came from. These are the
% standard loosened values: a float may take nine tenths of a page from the top,
% a page needs only seven percent of text to count as a text page, and a
% float-only page has to be three quarters full before LaTeX will make one.
\renewcommand{\topfraction}{0.90}
\renewcommand{\bottomfraction}{0.80}
\renewcommand{\textfraction}{0.07}
\renewcommand{\floatpagefraction}{0.75}
\setcounter{topnumber}{3}
\setcounter{bottomnumber}{2}
\setcounter{totalnumber}{4}
\usepackage{textcomp}
\usepackage{chngcntr}
% The collaboration styles its name with the word in italic and the B upright.
% One definition, so the whole thesis changes with it.
\newcommand{\qbounce}{\textit{q}B\textit{ounce}}

\counterwithin{figure}{section}
% Tables were left on a running counter, so they came out 2.1, 3.2, 4.3 --
% the section changed and the number did not restart.
\counterwithin{table}{section}
\counterwithin{equation}{section}
\renewcommand{\thetable}{\arabic{section}.\arabic{table}}
\usepackage{nicefrac}
\usepackage{braket}
\usepackage{comment}
\usepackage{array}
\usepackage{todonotes} % \todo{...} and \todo[inline]{...} for visible
                        % reminders -- see the "TODO (Loic)" boxes below
\PassOptionsToPackage{hyphens}{url}
\usepackage[linkbordercolor=TUmagenta,citebordercolor=TUgreen,
urlbordercolor=TUblue]{hyperref}
\usepackage[babel,german=guillemets]{csquotes}
% textgreek removed: nothing in the thesis uses it, and it pulls in the LGR
% encoding whose font descriptors (LGRcmr.fd) are absent from this TeX Live
% install -- which aborted the run with 'This NFSS system isn't set up
% properly'. Re-add it together with the cbfonts-fd package if Greek text
% (not math) is ever needed.
% \usepackage{textgreek}
% float provides the [H] placement used by two figures; without it LaTeX
% raises 'Unknown float option H' and drops the figure.
\usepackage{float}
% Architecture graphs live in ../architecture_graphs.tex and are \input by both
% this thesis and CodeArchitecture.tex, so the two can never drift apart.
\usepackage{pifont}   % \ding{172}.. circled numbers, matching the GUI's buttons
\usepackage{tikz}
\usetikzlibrary{arrows.meta,positioning,fit,backgrounds,calc}
\usepackage{enumerate}

\usepackage[most]{tcolorbox}

%%% For Bibliography
\usepackage[style=numeric-comp,backend=biber,defernumbers=true,
hyperref=true,sorting=none,giveninits=true,maxbibnames=99,
doi=true,isbn=false,url=false]{biblatex}
\AtEveryBibitem{\clearfield{month}} \AtEveryBibitem{\clearfield{day}}
\AtEveryBibitem{\clearfield{number}}

\bibliography{bibliography.bib}

%%% Header
\usepackage[automark]{scrlayer-scrpage}
\ihead{Bachelorarbeit}
\chead{}
\ohead{\Name}

%%% Margins, please play
\topmargin 0.0cm
\textheight 220mm                 % Height of text part of page
\textwidth=16cm
\oddsidemargin=0cm
\evensidemargin=0cm

% ---------------------------------------------------------------------------
% My \newcommands

\newcommand{\TUblue}{\color{TUblue}}
\newcommand{\TUmagenta}{\color{TUmagenta}}

\newcommand{\etal}{\textit{et al.}\xspace}
\newcommand{\figref}[1]{Figure~\ref{#1}}
\newcommand{\tabref}[1]{Table~\ref{#1}}

% Physics shorthands kept from the original draft
\newcommand{\UCN}{ultra-cold neutron}
\newcommand{\UCNs}{ultra-cold neutrons}
\newcommand{\GRS}{Gravity Resonance Spectroscopy}
% \peV, \neV and \mum used to live here.  They are gone on purpose: they were a
% second way of writing a unit next to \SI, and the micrometre ended up spelled
% \mum in one place and \,\mu\mathrm{m} in another.  Write \SI{1.41}{\pico
% \electronvolt} and \SI{5.87}{\micro\metre}.  Anything still using the old
% macros now fails to compile, which is the point.

\renewcommand{\labelenumii}{\arabic{enumi}.\arabic{enumii}}
\renewcommand{\labelenumiii}{\arabic{enumi}.\arabic{enumii}.\arabic{enumiii}}
\renewcommand{\labelenumiv}{\arabic{enumi}.\arabic{enumii}.\arabic{enumiii}.\arabic{enumiv}}

%%% Graphicspath
\graphicspath{{./}{titlepage/}}

% ---------------------------------------------------------------------------
% listings -- for the "Commented Code Excerpts" section
\usepackage{textcomp,listings}
\usepackage{lstautogobble}
\definecolor{codegreen}{rgb}{0,0.6,0}
\definecolor{codegray}{rgb}{0.5,0.5,0.5}
\definecolor{codepurple}{rgb}{0.58,0,0.82}
\definecolor{backcolour}{rgb}{0.95,0.95,0.92}
\definecolor{listgray}{rgb}{0.88,0.88,0.88}

\lstdefinestyle{mystyle}{
  xleftmargin=0.02\linewidth,
  xrightmargin=0.02\linewidth,
  framesep=1pt,
  backgroundcolor=\color{backcolour},
  commentstyle=\color{codegreen},
  keywordstyle=\color{magenta},
  numberstyle=\tiny\color{codegray},
  stringstyle=\color{codepurple},
  basicstyle=\ttfamily\footnotesize,
  columns=flexible,
  breakatwhitespace=false,
  breaklines=true,
  captionpos=b,
  keepspaces=true,
  numbers=left,
  numbersep=5pt,
  showspaces=false,
  showstringspaces=false,
  showtabs=false,
  frame=single,
  frameround=tttt,
  tabsize=2,
  autogobble=true,
  % The C++ and Python sources use em dashes and box-drawing rules inside their
  % comment banners. Under pdflatex with inputenc those bytes have no printable
  % meaning, and a single one inside a quoted range stops the build. Mapping
  % them to ASCII here means any range of the real sources can be quoted with
  % \lstinputlisting without having to edit the sources to suit the thesis.
  extendedchars=true,
  literate={—}{{-{}-{}-}}1 {─}{{-}}1 {–}{{-{}-}}1
}

\lstset{style=mystyle}

% Convenience languages for \lstinputlisting[language=...]
\lstdefinelanguage{None}{}
%-----------------------------------

% ---------------------------------------------------------------------------
% "Don't forget" box for the things Loic must fill in before submission.
% Usage:  \begin{remembertodo}{Short label}  ... \end{remembertodo}
\newtcolorbox{remembertodo}[1]{%
  colback=TUorange3, colframe=TUorange, colbacktitle=TUorange,
  coltitle=white, fonttitle=\bfseries\sffamily,
  title={TODO (Loïc) --- #1}, breakable, boxrule=0.8pt,
  left=6pt, right=6pt, top=4pt, bottom=4pt
}

%%% Local Variables:
%%% mode: latex
%%% TeX-master: "main.tex"
%%% End:
